\begin{abstract}
    LLM-based search agents achieve strong performance but suffer from severe latency, as each step requires serialized LLM reasoning followed by action of tool execution. 
    We revisit this bottleneck through the lens of speculation.
    While traditional predict-verify speculation paradigm can break serial execution, its benefit remains limited, as it retains the full original workload and adds extra inference overhead.
    We observe that early agent steps often involve simple evidence-gathering, where correct actions can often be predicted without full reasoning. 
    Building on these observations, we present \system, an algorithm-system co-design framework that expands the role of speculation in search agents to reduce latency.
    Algorithmically, \system introduces a two-phase adaptive speculation mechanism that selectively omits verification when safe.
    System-wise, a two-level scheduler regulates speculative requests based on engine load to ensure speculation remains beneficial.
    We implement \system in real-world systems. 
    Across extensive experimental settings, \system achieves up to $1.65\times$ end-to-end speedup while maintaining same or even achieving higher accuracy, enabling practical deployment of multi-step search agents.
\end{abstract}